\section{Limitations}
\label{sec:limitations}

The analytic results concern stated affine, local, or exact-memory assumptions; they do not predict a generic GRU or LSTM outcome. Their individual foundations are established, and the paper's contribution is their use as a coherent diagnostic contract.

The cue-memory audit has one training seed in each of four synthetic model/task cells, and none passes competence. It cannot estimate effects across training runs or establish repetition-induced forgetting. Probe regularization at tiny state scales also confounds comparisons across horizons. Float64 analysis of saved float32 states cannot restore lost numerical information.

The ordering intervention matches selected statistics, not every temporal dependency. Its final-motif exposure difference is structural (\cref{app:motif-exposure}); even a supported ordering effect would require further interventions to attribute it to one periodicity mechanism. Covariance dimension depends on conditioning, sample size, and coordinate metric; local tangent dynamics do not characterize global basin structure. State interventions can themselves cause off-manifold disruption. These limitations require explicit controls rather than a single collapse score.

The graph appendix covers one capped corrected GC-LSTM on one graph realization, with compact-output verification rather than independent binary-state reanalysis. It does not reproduce the complete original multi-model, five-seed figure or establish causal retention loss. No native cross-modal transfer is demonstrated; rendering one latent task in several modalities would not by itself establish such transfer. The proposed stronger competence and statistical design remains separate from the executed pilot evidence.
